\documentclass[11pt]{article}
\input{packages.tex}
\usepackage{ulem}
\DeclareMathOperator{\sign}{sign}
\DeclareMathOperator{\Imm}{Im}
\DeclareMathOperator{\DIF}{DIF}
\DeclareMathOperator{\rg}{rg}

% Настройка колонтитулов
\pagestyle{fancy}
\fancyhf{} % Очистить все поля
\fancyhead[L]{Как говорит наш любимый шеф, если человек \sout{техает матан} идиот, то это надолго}
\fancyhead[R]{\href{https://t.me/ceagest}{Автор}}
\fancyfoot[C]{\thepage}               % Справа внизу

\setlength{\headheight}{13.6pt}

\hypersetup{
    colorlinks=true,
    linkcolor=blue,
    citecolor=blue,
    filecolor=blue,
    urlcolor=blue,
}

\begin{document}

\begin{titlepage}
    \centering
    \vspace*{\fill}  % заполняет пространство сверху

    {\Huge Лекция 7. Теорема об открытом отображении. Основная геометрическая теорема. Теорема о производной обратного отображения.
    \par}

    \vspace*{\fill}  % заполняет пространство снизу
\end{titlepage}

\setcounter{page}{1}

\section{Теорема об открытом отображении}

\begin{lemma}
    Пусть $n \in \mathbb{N}$, $A \in \mathcal{L}\left(\mathbb{R}^n, \mathbb{R}^n\right)$. Пусть матрица линейного отображения $A$ в каноническом базисе обратима. Тогда
    \[\left\lVert A(x) - A(y) \right\rVert \geq \frac{1}{\left\lVert A^{-1} \right\rVert} \left\lVert x - y \right\rVert \quad \forall x, y \in \mathbb{R}^n\]
\end{lemma}

\begin{proof}
    Зафиксируем $x \in \mathbb{R}^n$. Так как линейное отображение $A$ обратимо, то по определению $x = A^{-1} (A(x))$. Тогда, 
    вспомнив лемму об оценке нормы значения линейного оператора через его норму и норму аргумента из конца пятой лекции, получаем:
    \[\left\lVert x \right\rVert = \left\lVert A^{-1} (A(x)) \right\rVert \leq \left\lVert A^{-1} \right\rVert \cdot \left\lVert A(x) \right\rVert\]
    При этом заметим, что $\left\lVert A^{-1} \right\rVert \neq 0$, так как иначе по свойствам нормы $A^{-1} \equiv 0$, что невозможно в силу отсутствия у тождественно нулевого отображения обратного.
    Тогда со спокойной душой делим обе части неравенства выше на эту норму:
    \[\forall x \in \mathbb{R}^n \hookrightarrow \left\lVert A(x) \right\rVert \geq \frac{\left\lVert x \right\rVert}{\left\lVert A^{-1} \right\rVert}\]
    Пользуясь полученным выше неравенством, уже совсем нетрудно показать требуемое. Достаточно лишь 
    вспомнить про линейность отображения $A$ и получить:
    \[\left\lVert A(x) - A(y) \right\rVert = \left\lVert A(x - y) \right\rVert \geq \frac{1}{\left\lVert A^{-1} \right\rVert} \cdot \left\lVert x - y \right\rVert \quad \forall x, y \in \mathbb{R}^n\]
\end{proof}

\begin{definition}
    Пусть $\Omega \subset \mathbb{R}^n$, $F \colon \Omega \to \mathbb{R}^n$. Пусть $x^0 \in \Omega$ --- внутренняя точка множества $\Omega$ и отображение $F$ дифференцируемо в точке $x^0$. Тогда \textbf{якобианом} отображения $F$ в точке $x^0$ будем называть 
    определитель матрицы Якоби отображения $F$ в точке $x^0$, то есть:
    \[\mathcal{J} F\left(x^0\right) = \det \mathcal{D} F\left(x^0\right)\]
\end{definition}

\begin{remark}
    Если кто-то однажды спросит вас в тёмной подворотне, можно ли рассматривать в определении якобиана отображение из подмножества $\mathbb{R}^n$ в $\mathbb{R}^m$, то есть брать в образе не обязательно равное по размерности арифметическое пространство,
    то смело отвечайте такому хитрецу, желающему одним вопросом отправить вас на пересдачу, что такое действо невозможно просто потому, что детерминант определён только для квадратных матриц, а значит размерности обязательно должны совпадать.
\end{remark}

\begin{theorem}[Основная геометрическая теорема]
    Пусть $\Omega \subset \mathbb{R}^n$, $F \colon \Omega \to \mathbb{R}^n$. Пусть $x^0 \in \Omega$ --- внутренняя точка множества $\Omega$ и отображение $F$ дифференцируемо в точке $x^0$, а также матрица Якоби $\mathcal{D} F\left(x^0\right)$ отображения $F$ в точке $x^0$ является обратимой (что равносильно не равенству нулю якобиана отображения $F$ в точке $x^0$). Тогда
    \[\exists \delta_0 > 0 \; \; \exists c > 0: \forall x \in \overline{B}_{\delta_0}\left(x^0\right) \hookrightarrow \left\lVert F\left(x\right) - F\left(x^0\right) \right\rVert \geq c \left\lVert x - x^0 \right\rVert\] 
\end{theorem}

\begin{proof}
    Так как $x^0$ --- внутренняя точка множества $\Omega$ и отображение $F$ дифференцируемо в точке $x^0$, то по определению:
    \[F\left(x\right) = F\left(x^0\right) + \underbrace{\mathcal{D} F \left(x^0\right) \left(x - x^0\right)}_{A\left(x - x^0\right)} + \varepsilon_{x^0} (x) \left\lVert x - x^0 \right\rVert \quad \forall x \in \overline{B}_{\delta}\left(x^0\right)\]
    Здесь, как и всегда, $\varepsilon_{x^0} (x) \to 0, \; x \to x^0$. При этом, если вдруг так получилось, что равенство из определения дифференцируемости выполнено для всех $x$ из открытого шара, то, уменьшив $\delta$ в два раза, получим выполнение равенства уже для нужного замкнутого шара.
    Далее, раз уж мы постарались и доказали предыдущую лемму, применим её для обратимого по условию линейного отображения, задающегося матрицей $\mathcal{D} F \left(x^0\right)$ (обозначим его $A$ для удобства):
    \begin{equation} \label{eq:geom:1}
        \left\lVert A \left(x - x^0\right) \right\rVert \geq \frac{1}{\left\lVert A^{-1} \right\rVert} \cdot \left\lVert x - x^0 \right\rVert
    \end{equation}
    Также в силу $\varepsilon_{x^0} (x) \to 0, \; x \to x^0$ получаем:
    \begin{equation} \label{eq:geom:2}
        \exists \widetilde{\delta} \in (0, \delta]: \forall x \in \overline{B}_{\widetilde{\delta}}\left(x^0\right) \hookrightarrow \left\lVert \varepsilon_{x^0} (x) \right\rVert \leq \frac{1}{2\left\lVert A^{-1} \right\rVert}
    \end{equation}
    Запомнив два неравенства выше, выпишем третье и приступим к оценке приращения. По неравенству треугольника для нормы:
    \begin{equation} \label{eq:geom:3}
        \left\lVert a \pm b \right\rVert \geq \left\lVert a \right\rVert - \left\lVert b \right\rVert
    \end{equation}
    Тогда, вспоминая равенство из определения дифференцируемости, оцениваем наконец искомую норму разности:
    \[\left\lVert F\left(x\right) - F\left(x^0\right) \right\rVert = \left\lVert A\left(x - x^0\right) + \varepsilon_{x^0} (x) \left\lVert x - x^0 \right\rVert \right\rVert \geq [\text{неравенство (\ref{eq:geom:3})}] \geq \]
    \[\geq \left\lVert A \left(x - x^0\right) \right\rVert - \left\lVert \varepsilon_{x^0} (x) \left\lVert x - x^0 \right\rVert \right\rVert \geq [\text{неравенство (\ref{eq:geom:1}) + свойство нормы}] \geq \]
    \[\geq \frac{1}{\left\lVert A^{-1} \right\rVert} \cdot \left\lVert x - x^0 \right\rVert - \left\lVert x - x^0 \right\rVert \cdot \left\lVert \varepsilon_{x^0} (x)\right\rVert = \left\lVert x - x^0 \right\rVert \left(\frac{1}{\left\lVert A^{-1} \right\rVert} - \left\lVert \varepsilon_{x^0} (x) \right\rVert \right) \geq \]
    \[\geq [\text{неравенство (\ref{eq:geom:2})}] \geq \left\lVert x - x^0 \right\rVert \cdot \frac{1}{2\left\lVert A^{-1} \right\rVert} \quad \forall x \in \overline{B}_{\widetilde{\delta}} \left(x^0\right)\] 
    Из вышеописанной цепочки неравенств нетрудно видеть, что, положив
    \[\delta_0 = \widetilde{\delta}, \quad c = \frac{1}{2\left\lVert A^{-1} \right\rVert},\]
    мы получим то, что и требовалось.
\end{proof}

\begin{theorem}[Об открытом отображении]
    Пусть $\Omega \subset \mathbb{R}^n$ --- открыто, $F \in \DIF (\Omega, \mathbb{R}^n)$. Пусть матрица Якоби $\mathcal{D} F (x)$ отображения $F$ обратима 
    для всех $x \in \Omega$ (то есть $\mathcal{J} F(x) \neq 0$ для всех $x \in \Omega$). Тогда множество $F(\Omega)$ открыто в $\mathbb{R}^n$.
\end{theorem}

\begin{remark}
    Условие не равенства нулю якобиана отображения $F$ всюду на множестве $\Omega$ важно и его, вообще говоря, нельзя отбросить.
    Действительно, приведём следующий пример:
    \[F(x) = x^2, \quad x \in (-1, 1)\]
    Понятно, что все требования на отображение выполнены, но при этом $\mathcal{D} F (0) = F'(0) = 0$ и множество $F(-1, 1) = [0, 1)$ не является открытым в $\mathbb{R}$.
\end{remark}

\begin{proof}
    Фиксируем $x^0 \in \Omega$ и $y^0 = F\left(x^0\right)$. Наша цель состоит в том, чтобы найти достаточно малый шар с центром в точке $y^0$, который целиком залезет в $F(\Omega)$, так как тогда по определению открытого множества в силу произвольного выбора точки $x^0$ получим требуемое.
    Воспользуемся же для достижения цели основной геометрической теоремой:
    \[\exists \delta > 0 \; \; \exists c > 0: \overline{B}_{\delta} \left(x^0\right) \subset \Omega \quad \text{ и } \quad \forall x \in \overline{B}_{\delta} \left(x^0\right) \hookrightarrow \left\lVert F \left(x\right) - F\left(x^0\right) \right\rVert \geq c \left\lVert x - x^0 \right\rVert\]
    Зафиксируем эти самые $\delta$ и $c$ из условия выше и покажем, что $B_{\frac{c \delta}{2}} \left(y^0\right) \subset F(\Omega)$. Возьмём для этого произвольно $y^{\ast} \in B_{\frac{c \delta}{2}} \left(y^0\right)$ и рассмотрим следующую функцию:
    \begin{figure}[h]
    \centering
    \includegraphics[width=0.7\linewidth]{pictures/Теорема об открытом отображении.jpg}
    \end{figure}
    \[h(x) = \left\lVert F(x) - y^{\ast} \right\rVert \]
    Сделаем для неё пару замечаний. Ясно, что значение $h$ в точке $x^0$ можно оценить так:
    \[h\left(x^0\right) = \left\lVert F\left(x^0\right) - y^{\ast}\right\rVert = \left\lVert y^0 - y^{\ast} \right\rVert < [\text{определение открытого шара}] < \frac{c \delta}{2}\]
    При этом для $h(x)$ при $x \in S_{\delta}\left(x^0\right) = \{x \in \mathbb{R}^n: \left\lVert x - x^0 \right\rVert = \delta\}$ можно выписать следующие неравенства:
    \[h(x) = \left\lVert F(x) - y^{\ast} \right\rVert \geq [\text{следствие из неравенства треугольника}] \geq\]
    \[\geq \left\lVert F(x) - y^0 \right\rVert - \left\lVert y^0 - y^{\ast} \right\rVert \geq [\text{основная геометрическая теорема}] \geq\]
    \[\geq \underbrace{c \left\lVert x - x^0 \right\rVert}_{= c \delta} - \left\lVert y^0 - y^{\ast} \right\rVert > [\text{определение открытого шара и сферы}] > c \delta - \frac{c \delta}{2} = \frac{c \delta}{2}\]
    Пробежав глазками цепочку неравенств, заключаем:
    \[\forall x \in S_{\delta} \left(x^0\right) \hookrightarrow h(x) > \frac{c \delta }{2}\]
    Заметим же теперь, что в силу дифференцируемости отображения $F$ всюду на $\Omega$, а значит и непрерывности этого отображения в каждой точке $\Omega$, имеем $h \in C \left( \overline{B}_{\delta}\left(x^0\right) \right)$. Но множество $\overline{B}_{\delta}\left(x^0\right)$, как легко видеть, является компактом,
    откуда по теореме Вейерштрасса функция $h$ достигает минимума на $\overline{B}_{\delta}\left(x^0\right)$. 
    \\
    К этому моменту внимательного читателя должен возникнуть естественный вопрос: зачем же мы доказывали все эти оценки на функцию $h$? На самом деле это было нужно 
    только для того, чтобы сейчас уверенно заявить: $h$ достигает своего минимума точно внутри шара $\overline{B}_{\delta}\left(x^0\right)$, то есть на множестве $B_{\delta} \left(x^0\right)$.
    Действительно, функция $h$ не может достигать минимума на границе шара $\overline{B}_{\delta}\left(x^0\right)$, так как она представляет собой сферу $S_{\delta} \left(x^0\right)$, все значения на которой по доказанному выше 
    заведомо больше, чем в центре шара.
    \\
    После ответа на предыдущий вопрос, читатель, скорее всего, находится в ещё большем замешательстве, ведь совсем непонятно, к чему нам понадобилось рассматривать именно внутренность замкнутого шара. Замысел прост: 
    мы хотим продифференцировать функцию $h$, а если быть точным, её квадрат, так как с обычной нормой каши не сваришь. Распишем для этого $h^2(x)$ по определению: 
    \[h^2(x) = \left\lVert F(x) - y^{\ast} \right\rVert^2 = \sum_{j = 1}^n \left(F_j (x) - y^{\ast}_j \right)^2\]
    Далее важно заметить, что если $x^{\ast}$ является точкой минимума функции $h$, то эта же точка 
    является точкой минимума для функции $h^2$. Кроме того, если $x^{\ast}$ --- точка минимума функции, то её вектор градиент в этой точке равен нулю. Тогда, расписывая равенство нулю градиента функции $h^2$ в точке $x^{\ast}$ (по третьему необходимому условию дифференцируемости этот вектор состоит из частных производных по всем аргументам рассматриваемой функции) и учитывая полученное выше выражение, имеем:
    \begin{equation*} \label{eq:th:1.2:1} \tag{$\ast$}
        \forall i \in \{1, \ldots, n\} \hookrightarrow \frac{\partial h^2(x)}{\partial x_i} \Bigg|_{x = x^{\ast}} = \sum_{j = 1}^n \frac{\partial \left(F_j(x) - y_j^{\ast} \right)^2}{\partial x_i} \Bigg|_{x = x^{\ast}} = 
    \sum_{j = 1}^n 2 \left(F_j(x^{\ast}) - y_j^{\ast} \right) \frac{\partial F_j}{\partial x_i} \left(x^{\ast}\right) = 0
    \end{equation*}
    Сохранив полученное в памяти, рассмотрим вектор $z$ и вспомним, из чего же состоит матрица Якоби отображения $F$ в точке $x^{\ast}$:
    \begin{equation*}
        z = 
        \begin{pmatrix}
            F_1 \left(x^{\ast} \right) - y_1^{\ast} \\
            \vdots \\
            F_n \left(x^{\ast} \right) - y_n^{\ast}
        \end{pmatrix}
        \quad \quad \quad 
        \mathcal{D} F \left(x^{\ast} \right) = 
        \begin{pmatrix}
            \displaystyle \frac{\partial F_1}{\partial x_1} \left(x^{\ast} \right) & \ldots & \displaystyle \frac{\partial F_1}{\partial x_n} \left(x^{\ast} \right) \\
            \displaystyle \vdots & \ddots & \vdots \\
            \displaystyle \frac{\partial F_n}{\partial x_1} \left(x^{\ast} \right) & \ldots & \displaystyle \frac{\partial F_n}{\partial x_n} \left(x^{\ast} \right)
        \end{pmatrix}
    \end{equation*}
    Методом пристального взгляда на два этих замечательных объекта, заключаем, что полученное таким трудом условие (\ref{eq:th:1.2:1}) равносильно следующему:
    \[\left(\mathcal{D} F \left(x^{\ast}\right) \right)^{T} z = 0\]
    Теперь осталось лишь заметить, что перед нами система линейных уравнений с равным числом уравнений и неизвестных, причём определитель матрицы этой системы не равен нулю, так как представляет собой 
    определитель матрицы Якоби отображения $F$ в точке $x^{\ast} \in \Omega$, который отличен от нуля в любой точке $\Omega$ по условию. Тогда, как известно из 
    курса линала, эта СЛУ имеет единственное решение. Но в силу однородности системы решение $z = 0$ есть всегда, а значит это и есть то самое единственное решение. Тогда, вспомнив, как же мы определяли вектор $z$, с радостью 
    заключаем, что 
    \[F \left(x^{\ast} \right) - y^{\ast} = 0 \iff F \left(x^{\ast} \right) = y^{\ast}\]
    Из полученного выше ясно, что мы нашли прообраз для $y^{\ast}$, а значит $y^{\ast} \in F(\Omega)$. Но $y^{\ast} \in B_{\frac{c \delta}{2}} \left(y^0\right)$ было выбрано произвольно, откуда $B_{\frac{c \delta}{2}} \left(y^0\right) \subset F(\Omega)$. Теперь уж внимательный читатель сразу заметит, что это победа!
\end{proof}

\section{Теорема о производной обратного отображения}

\begin{theorem}[О производной обратного отображения]
    Пусть $\Omega \subset \mathbb{R}^n$, $x^0$ --- внутренняя точка множества $\Omega$.
    Пусть:
    \begin{enumerate}
        \item Отображение $F \colon \Omega \to \mathbb{R}^n$ обратимо на $\Omega$ (то есть $\exists F^{-1} \colon F(\Omega) \to \Omega$)
        \item Отображение $F$ дифференцируемо в точке $x^0$
        \item $\mathcal{J} F\left(x^0\right) \neq 0$
        \item Точка $y^0 = F\left(x^0\right)$ является внутренней точкой множества $F(\Omega)$
        \item Обратное отображение $F^{-1}$ непрерывно в точке $y^0$
    \end{enumerate}
    Тогда отображение $F^{-1}$ дифференцируемо в точке $y^0$ и, более того, справедливо равенство:
    \[\mathcal{D} F^{-1} \left(y^0\right) = \left(\mathcal{D} F \left(x^0\right) \right)^{-1}\]
\end{theorem}

\begin{proof}
    Так как отображение $F$ дифференцируемо в точке $x^0$, являющейся внутренней точкой множества $\Omega$, то 
    по определению дифференцируемости:
    \begin{equation*} \label{eq:th:2.1:1} \tag{$\star$}
        \exists \delta > 0: F\left(x\right) - F\left(x^0\right) = \mathcal{D} F \left(x^0\right) \left(x - x^0\right) + \varepsilon_{x^0} (x) \left\lVert x - x^0 \right\rVert \quad \forall x \in B_{\delta} \left(x^0\right)
    \end{equation*}
    Здесь $\varepsilon_{x^0} (x) \to 0, \; x \to x^0$. Далее в силу существования обратного отображения $F^{-1}$ и его непрерывности в точке $y^0$ получаем:
    \[\exists \varepsilon > 0: \forall y \in B_{\varepsilon} \left(y^0\right) \hookrightarrow F^{-1} (y) \in B_{\delta} \left(x^0\right)\]
    Коль так, то для всех $y \in B_{\varepsilon}\left(y^0\right)$ можем подставить в условие (\ref{eq:th:2.1:1}) значения 
    \[x = F^{-1}(y) \in B_{\delta}\left(x^0\right) \quad \text{ и } \quad x^0 = F^{-1}\left(y^0\right)\]
    Давайте же сделаем это:
    \[y - y^0 = \mathcal{D} F \left(F^{-1}\left(y^0\right)\right) \left(F^{-1}\left(y\right) - F^{-1}\left(y^0\right)\right) + \varepsilon_{F^{-1}\left(y^0\right)} \left(F^{-1}\left(y\right)\right) \left\lVert F^{-1}\left(y\right) - F^{-1}\left(y^0\right) \right\rVert \]
    Так как $\mathcal{J} F \left(x^0\right) \neq 0$, то матрица Якоби $\mathcal{D} F \left(x^0\right) = \mathcal{D} F\left(F^{-1}\left(y^0\right)\right)$ обратима. Умножим тогда обе части равенства выше на обратную матрицу к этой матрице Якоби:
    \[\left(\mathcal{D} F \left(F^{-1}\left(y^0\right)\right)\right)^{-1} \left(y - y^0\right) = F^{-1}\left(y\right) - F^{-1}\left(y^0\right) + \]
    \[+ \left(\mathcal{D} F \left(F^{-1}\left(y^0\right)\right)\right)^{-1} \varepsilon_{F^{-1}\left(y^0\right)} \left(F^{-1}\left(y\right)\right) \left\lVert F^{-1}\left(y\right) - F^{-1}\left(y^0\right) \right\rVert\]
    Чтобы не утонуть в выкладках, обозначим линейное отображение $\left(\mathcal{D} F \left(F^{-1}\left(y^0\right)\right)\right)^{-1}$ за $A$ (линейное отображение однозначно задаётся своей матрицей) и перекинем второе слагаемое в правой части равенства налево:
    \[F^{-1}\left(y\right) - F^{-1}\left(y^0\right) = A\left(y - y^0\right) - A \cdot \varepsilon_{F^{-1}\left(y^0\right)} \left(F^{-1}\left(y\right)\right) \left\lVert F^{-1}\left(y\right) - F^{-1}\left(y^0\right) \right\rVert \quad  \forall y \in B_{\varepsilon}\left(y^0\right)\]
    Смекалистый читатель приметит, что это равенство уж очень похоже на определение дифференцируемости отображения $F^{-1}$ в точке $y^0$. Не хватает только тайного знания 
    о том, что второе слагаемое в правой части равенства есть $\overline{o} \left(\left\lVert y - y^0 \right\rVert \right), \; y \to y^0$. Для того, чтобы это увидеть, воспользуемся основной геометрической теоремой для отображения $F$ и точки $x^0$ (понятно, что все условия для её использования выполнены):
    \[\left\lVert y - y^0 \right\rVert = \left\lVert F\left(x\right) - F\left(x^0\right) \right\rVert \geq c \left\lVert x - x^0 \right\rVert = c \left\lVert F^{-1}\left(y\right) - F^{-1}\left(y^0\right)\right\rVert \quad \forall x \in \overline{B}_{\delta}\left(x^0\right)\]
    \[\iff\]
    \begin{equation*} \label{eq:th:2.1:2} \tag{$\star \star$}
        \left\lVert F^{-1}\left(y\right) - F^{-1}\left(y^0\right)\right\rVert \leq \frac{1}{c} \left\lVert y - y^0 \right\rVert \quad \forall y \in B_{\varepsilon}\left(y^0\right)
    \end{equation*}
    Важно отметить, что здесь $\delta$ выбирается, вообще говоря, независимо по отношению к выбору $\delta$ в условии (\ref{eq:th:2.1:1}), а значит, если справедливо (\ref{eq:th:2.1:1}), то не обязательно 
    неравенство из основной геометрической теоремы справедливо в замкнутом шаре этого же радиуса. Но мы не переживаем по этому поводу, и говорим волшебные слова о том, что, если потребуется, 
    мы подожмём $\delta$ в двух этих условиях и обозначим полученное вновь за $\delta$, и тогда оба условия будут справедливы в соответствующих шарах этого радиуса. 
    \\
    От неравенства (\ref{eq:th:2.1:2}) мы возьмём только ограниченность функции
    $\frac{\left\lVert F^{-1}\left(y\right) - F^{-1}\left(y^0\right)\right\rVert}{\left\lVert y - y^0 \right\rVert}$
    на множестве $B_{\varepsilon}\left(y^0\right)$, а далее отметим, что следующее условие выполнено по теореме о пределе композиции отображений:
    \[A \cdot \varepsilon_{F^{-1}\left(y^0\right)} \left(F^{-1}(y)\right) \to 0, \; y \to y^0\]
    В самом деле, $F^{-1}(y) \to F^{-1}\left(y^0\right), \; y \to y^0$ из непрерывности обратного отображения в точке $y^0$, и при этом $A$ --- некоторая фиксированная матрица.
    \\
    Осталось сделать лишь пару штрихов, первым из которых будет такое обозначение:
    \[\widetilde{\varepsilon}_{y^0} (y) := \underbrace{-A \cdot \varepsilon_{F^{-1}\left(y^0\right)} \left(F^{-1}(y)\right)}_{\text{б.м. при } y \to y^0} \cdot \underbrace{\frac{\left\lVert F^{-1}\left(y\right) - F^{-1}\left(y^0\right)\right\rVert}{\left\lVert y - y^0 \right\rVert}}_{\text{огр. при } y \to y^0}\]
    Нетрудно видеть из всех наших потугов, что $\widetilde{\varepsilon}_{y^0} (y) \to 0, \; y \to y^0$. Вторым же штрихом будет причёсывание того самого второго слагаемого в равенстве для приращения обратного отображения:
    \[-A \cdot \varepsilon_{F^{-1}\left(y^0\right)} \left(F^{-1}\left(y\right)\right) \left\lVert F^{-1}\left(y\right) - F^{-1}\left(y^0\right) \right\rVert = \]
    \[= -A \cdot \varepsilon_{F^{-1}\left(y^0\right)} \left(F^{-1}\left(y\right)\right) \frac{\left\lVert F^{-1}\left(y\right) - F^{-1}\left(y^0\right) \right\rVert}{\left\lVert y - y^0 \right\rVert} \left\lVert y - y^0 \right\rVert =\]
    \[= \widetilde{\varepsilon}_{y^0} (y) \left\lVert y - y^0 \right\rVert = \overline{o} \left(\left\lVert y - y^0 \right\rVert\right), \; y \to y^0\]
    Всё, теперь осталось только подставить вместо этого слагаемого полученное представление:
    \[F^{-1}\left(y\right) - F^{-1}\left(y^0\right) = A\left(y - y^0\right) + \overline{o} \left(\left\lVert y - y^0 \right\rVert\right), \; y \to y^0\]
    Ура, перед нами заветное определение дифференцируемости отображения $F^{-1}$ в точке $y^0$. Более того, из этого же условия вытекает 
    требуемое равенство для матриц Якоби:
    \[\left(\mathcal{D} F \left(x^0\right) \right)^{-1} = A = \mathcal{D} F^{-1} \left(y^0\right)\]
    Теперь уже со спокойной душой заключаем, что требуемое доказано.
\end{proof}

\end{document}